\chapter{Спектральная счётная мера четырёхмерного объёма ячейки}
\label{ch:v10-cell-birth-four-volume-spectral-counting}

\section{Точная конечная спектральная модель}

Проверим, может ли размерность старой столкновительной ячейки $K_{43}$
выбрать её собственную длину. Возьмём безразмерный оператор формы
\begin{equation}
 \widehat D_{43}=\operatorname{diag}(0,1,\ldots,42),
 \qquad D_{\rm cell}=\frac{\widehat D_{43}}{\ell_{\rm cell}}.
 \label{eq:v10-spectral-counting-operator}
\end{equation}
При верхнем пороге $\Lambda=42/\ell_{\rm cell}$ счётный проектор имеет
ранг $43$. Дополнительно
\begin{equation}
 \ell_{\rm cell}^2\operatorname{Tr}D_{\rm cell}^2
 =\sum_{n=0}^{42}n^2=25585.
 \label{eq:v10-spectral-counting-moment}
\end{equation}
Таким образом, спектральная форма и кратность определены точно.

\section{Что выбирает счётная мера}

Верхний уровень фиксирует лишь произведение
\begin{equation}
 \Lambda\ell_{\rm cell}=42.
 \label{eq:v10-spectral-counting-product}
\end{equation}
Преобразование
\begin{equation}
 \Lambda\mapsto\frac{\Lambda}{s},
 \qquad \ell_{\rm cell}\mapsto s\ell_{\rm cell}
 \label{eq:v10-spectral-counting-orbit}
\end{equation}
сохраняет все уровни относительно порога, число состояний и нормированные
спектральные моменты.

В логарифмических координатах родитель счётного условия имеет вид
\begin{equation}
 \mathcal P_{\rm count}
 =\frac12\bigl(\log\Lambda+\log\ell_{\rm cell}-\log42\bigr)^2.
 \label{eq:v10-spectral-counting-parent}
\end{equation}
Его гессиан равен
\begin{equation}
 \begin{pmatrix}1&1\\1&1\end{pmatrix},
 \qquad \rank=1,
 \qquad \dim\ker=1,
 \qquad \operatorname{Spec}=\{0,2\}.
 \label{eq:v10-spectral-counting-hessian}
\end{equation}
Нулевая мода $(-1,1)^T$ совпадает с орбитой обрезания и длины.

\begin{theorem}[Спектральный счёт выбирает только относительную меру]
Конечный спектр из $43$ уровней точно определяет счёт, верхний порог и
нормированные моменты, но фиксирует только $\Lambda\ell_{\rm cell}$. Поэтому
он не выбирает отдельно $\ell_{\rm cell}$ и
$v_{\rm cell}=\ell_{\rm cell}^4$.
\end{theorem}

\begin{nogocriterion}[Запрет превращения размерности в объём]
Число внутренних состояний не является физическим четырёхмерным объёмом.
Без независимо выбранного спектрального обрезания равенство
$\Lambda\ell_{\rm cell}=42$ сохраняет непрерывную масштабную орбиту.
\end{nogocriterion}

\begin{protocol}[Статус спектральной счётной меры]\begin{enumerate}
\item конечный спектральный носитель: \textbf{$43/43$};
\item точная счётная мера и момент: \textbf{получены};
\item архитектура: \textbf{$8/8$};
\item происхождение относительной меры: \textbf{$3/5$};
\item ранг и размерность ядра: \textbf{$1/1$};
\item абсолютный объём ячейки: \textbf{$0/1$};
\item абсолютная энергия часов: \textbf{$0/1$};
\item следующий гейт: \textbf{аудит топологических квантов объёма}.
\end{enumerate}\end{protocol}

Машинный сертификат сохранён в
\path{s2t_v10_cell_birth_four_volume_spectral_counting_measure_origin_gate_results.json}.